AviSynth "--- это скриптовый фреймсервер: файл с расширением \texttt{.avs}
описывает цепочку фильтров, применяемых к видеопотоку, а результат
передаётся любой программе, умеющей открывать видео~\cite{avisynth}.
Собственного графического интерфейса AviSynth не имеет: его роль заключается
в том, чтобы вычислять кадры по требованию, не создавая промежуточных
файлов.

В качестве внешней программы выступает VirtualDub2: он открывает скрипт как
обычный видеофайл, позволяет просмотреть результат и выполнить финальное
кодирование~\cite{virtualdub}. Такая связка покрывает типичные задачи
монтажа домашнего видео: вырезание фрагментов, обрезку краёв кадра,
изменение разрешения и частоты кадров.

Ниже рассматриваются модель фреймсервера, базовые команды AviSynth
и соотношения, связывающие параметры монтажа со свойствами видеопотока,
а также приводится пример конвейера обработки.
